\chapter{The Agent Layer: MCP, and What's Next}

\section{Why an MCP server on top of a CLI}

The CLI (\texttt{qp demo}, \texttt{qp backtest}, \texttt{qp walkforward}, \texttt{qp
paper}) is sufficient for a human at a terminal. An agent needs the same
capabilities as typed, discoverable tool calls rather than a stream of shell
output to parse -- and once the pipeline speaks MCP, it is reachable from
Claude Code, Claude Desktop, or any other MCP client (including, per the brief
this project started from, an orchestration layer like Google Stitch) without
writing a second integration.

\section{What is exposed today}

\begin{table}[h]
\footnotesize
\centering
\begin{tabularx}{\textwidth}{L{4.6cm} X L{2.6cm}}
\toprule
\thh{Tool} & \thh{Effect} & \thh{Side effects} \\
\midrule
\texttt{describe\_pipeline()} & returns this project's README & none \\
\texttt{list\_configs()} & lists \texttt{configs/*.yaml} & none \\
\texttt{list\_runs()} & lists \texttt{runs/*} with Sharpe/CAGR/maxDD & none \\
\texttt{run\_backtest(config\_path)} & full data$\to$features$\to$signal$\to$backtest$\to$tearsheet & writes \texttt{runs/<name>/} \\
\texttt{get\_run\_report(name)} & plain-text metrics table & none \\
\texttt{get\_tearsheet\_path(name)} & absolute path to the PNG & none \\
\texttt{paper\_iteration(config\_path)} & one live-loop step against the paper broker & writes \texttt{execution\_state/} \\
\bottomrule
\end{tabularx}
\caption{\texttt{mcp/server.py}, built on \texttt{FastMCP}. Registered for Claude
Code via the project's \texttt{.mcp.json}. Nothing here can reach a real broker
(Chapter~7).}
\end{table}

The server degrades gracefully: without the \texttt{mcp} package installed, it
prints install instructions to stderr and exits 0 rather than raising inside an
agent's tool-discovery call.

\section{Roadmap}

\begin{enumerate}
  \item \textbf{Now.} Read-mostly tools above; an agent proposes configs and
  reads results.
  \item \textbf{Next -- \texttt{propose\_signal(description)}.} The calling agent
  drafts a new \texttt{Signal} subclass from a natural-language description,
  drops it into \texttt{src/qp/signals/}, registers it, and backtests it against
  the existing baselines in one call.
  \item \textbf{Next -- \texttt{sweep(param\_grid)}.} A walk-forward parameter
  study across a grid (e.g.\ \texttt{lookback} $\in \{20, 60, 120\}$), returning
  a ranked table -- the mechanical half of strategy research, delegated.
  \item \textbf{Next -- \texttt{research\_report(run\_a, run\_b)}.} Diff two runs'
  metrics, weights, and cost decomposition into a short markdown memo.
  \item \textbf{Later -- real data \& a real signal.} Swap the \texttt{csv}
  source for an ENTSO-E day-ahead export and add a power-market
  \texttt{Signal} (features: load, wind/solar, and fuel-price forecasts; target:
  the day-ahead$\leftrightarrow$intraday spread, or a storage-dispatch P\&L).
  This reuses the PyPSA / OR-Tools background directly and is the artefact
  described in Chapter~10 as the actual point of this project.
\end{enumerate}
